\section{System overview and management model}
\label{sec:system}

The parsing service recovers document content, the web application
maintains collections and their files, and the MCP server exposes these
operations to agents \cite{mcp}. Collection placement defines where a
document is visible, its description guides selection within the
collection, and
its index entries make passages available to lexical and dense search.
Reconciliation checks for selected management operations compare the
document files, tree registry and per-library metadata.

\subsection{Management levels and the tiered access path}
\label{sec:levels}

Collections define ownership and scope; documents contain text and
descriptive metadata. A document may be stored in several parts, which the
index splits into segments called chunks. At query time, a passage is any
text span returned or scored. The catalogue lists collections, a shelf lists
documents within one collection, the document view exposes the body, and the
segment view exposes passages with source references. Table~\ref{tab:levels}
maps each level to the operations that maintain and check it.

In the implementation, a collection is called a library or knowledge base;
this paper uses \emph{collection} for the managed object and retains
\emph{library} only in implementation terms such as ``sub-library'' and
``library registry.'' The evidence gate is the Laya passage scorer; an
unqualified \emph{judge} in the evaluation refers to the blind answer
evaluator, not the gate.

\begin{table*}[t]
\centering
\small
\setlength{\tabcolsep}{4pt}
\renewcommand{\arraystretch}{1.15}
\caption{\looseness=-1 Management levels and the tiered access path. Each level names the
objects it maintains, the operations that write them, the state that describes
them, the tier that reads them, and the measurement that exercises it.
Evidence labels E1--E7 denote the evaluation campaigns of
Section~\ref{sec:eval}. The selected access policy determines which tiers are
visited; a pinned search can bypass catalogue navigation.}
\label{tab:levels}
\begin{tabular}{@{}>{\raggedright\arraybackslash}p{0.12\textwidth}>{\raggedright\arraybackslash}p{0.21\textwidth}>{\raggedright\arraybackslash}p{0.18\textwidth}>{\raggedright\arraybackslash}p{0.12\textwidth}>{\raggedright\arraybackslash}p{0.12\textwidth}>{\raggedright\arraybackslash}p{0.13\textwidth}@{}}
\toprule
Level & Objects and operations & State description & Read tier & Check & Evidence\\
\midrule
Collection & libraries and sub-libraries; agent assignment from parsed content; description and tag derivation; catalogue listing & library registry, per-library metadata, library-level index & catalogue & catalogue inventory and description round-trip & E1, E3\\
Document & documents and their parts; structural split; duplicate review; move, rename, delete; forced reindex & document files, tree registry, per-library metadata & shelf and document & reconciliation on selected operations & E1\\
Segment & dense and lexical indices; optional evidence gate & chunk identifiers, per-segment scores with provenance & segment & index coverage and post-index probe & E1, E2, E4, E6, E7\\
\bottomrule
\end{tabular}
\end{table*}

The ingest agent uses catalogue descriptions to place parsed material and
derives document descriptions and tags from its content. Sub-libraries are
tree nodes with their own descriptions and indices. At the document level,
duplicate review identifies candidates and a lifecycle operation can move,
rename or delete a document. At the segment level, search returns indexed
chunks, while the optional evidence gate fetches candidate text and scores
query-conditioned passages with provenance.

\looseness=-1 Figure~\ref{fig:levels} illustrates how catalogue-first search narrows
from collections to passages. Pinned search can reach indexed chunks
directly within specified collections. Inspecting the catalogue and shelf
separately helps locate a missing collection or document.

\begin{figure}[t]
\centering
\includegraphics[width=\columnwidth]{fig_library_analogy_vector}
\caption{\looseness=-1 Library analogy for a catalogue-first HMRS access path. The library catalogue, shelf, book and passage correspond to the collection catalogue, document listing, document body and segment. Evidence scoring is optional; pinned search may bypass the catalogue and shelf.}
\label{fig:levels}
\end{figure}

\setcounter{dbltopnumber}{2}
\begin{figure*}[t]
\centering
\includegraphics[width=\textwidth]{fig_service_architecture_vector}
\caption{HMRS service architecture. The browser and agent/CLI access the web and MCP interfaces. Parsing, management and retrieval operate on shared document files, the tree registry, library metadata and search indices. Laya scoring is optional; the Neo4j graph supports navigation and optional candidate expansion.}
\label{fig:arch}
\end{figure*}

Figure~\ref{fig:arch} shows which services write and expose the shared
state. Table~\ref{tab:functions} then identifies where each management or
retrieval operation is checked. The ingest agent chooses a collection,
the save service validates the destination identifier, and an audit checks
the result that was stored.

\begin{table*}[t]
\centering
\small
\setlength{\tabcolsep}{4pt}
\caption{Management functions, the enforcement or workflow check associated
with each, and the evaluation that exercises it. The table separates API/MCP
controls from agent procedures and campaign audits; semantic choices such as
which collection owns a document are made and recorded by the workflow.
Evidence labels E1--E7 denote the campaigns in Section~\ref{sec:eval}.}
\label{tab:functions}
\begin{tabular}{@{}>{\raggedright\arraybackslash}p{0.13\textwidth}>{\raggedright\arraybackslash}p{0.37\textwidth}>{\raggedright\arraybackslash}p{0.30\textwidth}>{\raggedright\arraybackslash}p{0.11\textwidth}@{}}
\toprule
Group & Operations & Platform or workflow check & Evidence\\
\midrule
Acquisition & submit, parse (MinerU), quality review, duplicate review, structural split, store, index & parse-task state; workflow quality checks; entry-point-specific deduplication and indexing; post-index probe & E1\\
Organisation & create library or sub-library, content-informed placement, descriptions and tags, duplicate review & API validates identifiers and required fields; workflow checks descriptions and tags; audit compares files, tree and library metadata & E1\\
Lifecycle & move, rename, delete, force reindex, graph build & operation-specific API updates; reconciliation audit detects selected stale or inconsistent state & E1\\
Query & tier selection, vector, keyword, catalogue and two-stage search, gate invocation & endpoint enforces tool allow/deny lists, timeouts and turn limits; invoked gate returns survivor provenance & E2, E3, E4, E6, E7\\
Experience & extract, approve, tier, search & experience tools apply their own retrieval, verification and credibility filters & E5\\
\bottomrule
\end{tabular}
\end{table*}

\begin{figure*}[t]
\centering
\includegraphics[width=\textwidth]{fig_workflow_vector}
\caption{\looseness=-1 Illustrative HMRS workflow from document acquisition to evidence-based reporting. Documents are parsed when needed, organised into collections and indexed according to the entry point. A question can use catalogue navigation or pinned search; optional Laya scoring precedes evidence inspection. A completed review without support is distinct from an incomplete search or tool failure, which may occur at any stage.}
\label{fig:workflow}
\end{figure*}

\subsection{Data model and object lifecycle}
\label{sec:datamodel}
\label{sec:lifecycle}

Documents are stored as Markdown with associated assets. A managed
collection's state spans the document files, \texttt{.tree-fs.json} and the
library's \texttt{.knowledge-base.yml}. The ingest workflow distinguishes
parsed, stored and indexed content and records completion after a retrieval
probe. The MCP save path can start background indexing; the web parsed-file
save path exposes indexing as a separate operation. After a move, rename
or reindex, the file path and metadata show the stored state and a
retrieval probe checks which path the index returns.
Sections~\ref{sec:acquisition}--\ref{sec:implementation} describe these
layers in turn: acquisition, organisation and lifecycle, tiered access
and the evidence gate, and the implementation.
